\documentclass{amsart}
% For dealing with references we use the comment environment
\usepackage{verbatim}
\newenvironment{reference}{\comment}{\endcomment}
%\newenvironment{reference}{}{}
\newenvironment{slogan}{\comment}{\endcomment}
\newenvironment{history}{\comment}{\endcomment}

% For commutative diagrams we use Xy-pic
\usepackage[all]{xy}

% We use 2cell for 2-commutative diagrams.
\xyoption{2cell}
\UseAllTwocells

% We use multicol for the list of chapters between chapters
\usepackage{multicol}

% This is generally recommended for better output
\usepackage{lmodern}
\usepackage[T1]{fontenc}

% For cross-file-references

% Package for hypertext links:
\usepackage{hyperref}

% For any local file, say "hello.tex" you want to link to please

% Theorem environments.
%
\theoremstyle{plain}
\newtheorem{theorem}[subsection]{Theorem}
\newtheorem{proposition}[subsection]{Proposition}
\newtheorem{lemma}[subsection]{Lemma}

\theoremstyle{definition}
\newtheorem{definition}[subsection]{Definition}
\newtheorem{example}[subsection]{Example}
\newtheorem{exercise}[subsection]{Exercise}
\newtheorem{situation}[subsection]{Situation}

\theoremstyle{remark}
\newtheorem{remark}[subsection]{Remark}
\newtheorem{remarks}[subsection]{Remarks}

\numberwithin{equation}{subsection}

% Macros
%
\def\lim{\mathop{\mathrm{lim}}\nolimits}
\def\colim{\mathop{\mathrm{colim}}\nolimits}
\def\Spec{\mathop{\mathrm{Spec}}}
\def\Hom{\mathop{\mathrm{Hom}}\nolimits}
\def\Ext{\mathop{\mathrm{Ext}}\nolimits}
\def\SheafHom{\mathop{\mathcal{H}\!\mathit{om}}\nolimits}
\def\SheafExt{\mathop{\mathcal{E}\!\mathit{xt}}\nolimits}
\def\Sch{\mathit{Sch}}
\def\Mor{\mathop{\mathrm{Mor}}\nolimits}
\def\Ob{\mathop{\mathrm{Ob}}\nolimits}
\def\Sh{\mathop{\mathit{Sh}}\nolimits}
\def\NL{\mathop{N\!L}\nolimits}
\def\CH{\mathop{\mathrm{CH}}\nolimits}
\def\proetale{{pro\text{-}\acute{e}tale}}
\def\etale{{\acute{e}tale}}
\def\QCoh{\mathit{QCoh}}
\def\Ker{\mathop{\mathrm{Ker}}}
\def\Im{\mathop{\mathrm{Im}}}
\def\Coker{\mathop{\mathrm{Coker}}}
\def\Coim{\mathop{\mathrm{Coim}}}

% Boxtimes
%
\DeclareMathSymbol{\boxtimes}{\mathbin}{AMSa}{"02}

%
% Macros for moduli stacks/spaces
%
\def\QCohstack{\mathcal{QC}\!\mathit{oh}}
\def\Cohstack{\mathcal{C}\!\mathit{oh}}
\def\Spacesstack{\mathcal{S}\!\mathit{paces}}
\def\Quotfunctor{\mathrm{Quot}}
\def\Hilbfunctor{\mathrm{Hilb}}
\def\Curvesstack{\mathcal{C}\!\mathit{urves}}
\def\Polarizedstack{\mathcal{P}\!\mathit{olarized}}
\def\Complexesstack{\mathcal{C}\!\mathit{omplexes}}
% \Pic is the operator that assigns to X its picard group, usage \Pic(X)
% \Picardstack_{X/B} denotes the Picard stack of X over B
% \Picardfunctor_{X/B} denotes the Picard functor of X over B
\def\Pic{\mathop{\mathrm{Pic}}\nolimits}
\def\Picardstack{\mathcal{P}\!\mathit{ic}}
\def\Picardfunctor{\mathrm{Pic}}
\def\Deformationcategory{\mathcal{D}\!\mathit{ef}}

\setlength{\emergencystretch}{3em}
\begin{document}
\title[Stacks Project, Tag 0EA4]{326 --- Purity of the ramification locus for finite-type maps from normal schemes to regular schemes}
\maketitle
\noindent Copyright (C) 2005--2025 Johan de Jong. Stacks Project authors. Source: \href{https://stacks.math.columbia.edu/tag/0EA4}{Tag 0EA4}.

\noindent Editorial difficulty note (not author text). Difficulty H3: Unlike purity for already quasi-finite maps, this statement must first force quasi-finiteness from codimension-one etaleness. The proof uses completion and normalized hypersurface sections to show that the exceptional fibre point is closed, then reduces the remaining ramification to branch-locus purity..

\noindent Editorial provenance note (not author text). History: excerpt from revision \texttt{a0444\allowbreak 6e57e\allowbreak c1fbc\allowbreak 252a8\allowbreak 71afc\allowbreak ec775\allowbreak 2fb28\allowbreak 07b14}, pione.tex, lines 6307--6383. Reference presentation was adapted; original-author context and core support blocks were retained or added. The mathematical wording is unchanged. Licensed under GNU FDL 1.2 or later; no invariant sections or cover texts. See ../licenses/COPYING. Context blocks reproduce author definitions and supporting results; they are not additional counted problems.

\begin{definition}
\label{morphisms-definition-unramified}
Let $f : X \to S$ be a morphism of schemes.
\begin{enumerate}
\item We say that $f$ is {\it unramified at $x \in X$} if
there exists an affine open neighbourhood $\Spec(A) = U \subset X$
of $x$ and affine open $\Spec(R) = V \subset S$
with $f(U) \subset V$ such that the induced ring map
$R \to A$ is unramified.
\item We say that $f$ is {\it G-unramified at $x \in X$} if
there exists an affine open neighbourhood $\Spec(A) = U \subset X$
of $x$ and affine open $\Spec(R) = V \subset S$
with $f(U) \subset V$ such that the induced ring map
$R \to A$ is G-unramified.
\item We say that $f$ is {\it unramified} if it is unramified
at every point of $X$.
\item We say that $f$ is {\it G-unramified} if it is G-unramified
at every point of $X$.
\end{enumerate}
\end{definition}

\begin{definition}
\label{morphisms-definition-etale}
Let $f : X \to S$ be a morphism of schemes.
\begin{enumerate}
\item We say that $f$ is {\it \'etale at $x \in X$} if
there exists an affine open neighbourhood $\Spec(A) = U \subset X$
of $x$ and affine open $\Spec(R) = V \subset S$
with $f(U) \subset V$ such that the induced ring map
$R \to A$ is \'etale.
\item We say that $f$ is {\it \'etale} if it is \'etale at every point of $X$.
\item A morphism of affine schemes $f : X \to S$ is called
{\it standard \'etale} if $X \to S$ is isomorphic to
$$
\Spec(R[x]_h/(g)) \to \Spec(R)
$$
where $R \to R[x]_h/(g)$ is a standard \'etale ring map, see
Algebra, Definition \ref{algebra-definition-standard-etale},
i.e., $g$ is monic and $g'$ invertible in $R[x]_h/(g)$.
\end{enumerate}
\end{definition}

\begin{definition}
\label{properties-definition-Cohen-Macaulay}
Let $X$ be a scheme. We say $X$ is {\it Cohen-Macaulay} if
for every $x \in X$ there exists an affine open neighbourhood
$U \subset X$ of $x$ such that the ring $\mathcal{O}_X(U)$ is
Noetherian and Cohen-Macaulay.
\end{definition}

\begin{definition}
\label{more-morphisms-definition-lci}
Let $f : X \to S$ be a morphism of schemes.
\begin{enumerate}
\item Let $x \in X$. We say that $f$ is {\it Koszul at $x$} if $f$
is of finite type at $x$ and there exists an open neighbourhood
and a factorization of $f|_U$ as $\pi \circ i$ where $i : U \to P$
is a Koszul-regular immersion and $\pi : P \to S$ is smooth.
\item We say $f$ is a {\it Koszul morphism}, or that
$f$ is a {\it local complete intersection morphism}
if $f$ is Koszul at every point.
\end{enumerate}
\end{definition}

\begin{situation}
\label{limits-situation-descent-property}
Let $S = \lim S_i$ be a limit of a directed system of schemes
with affine transition morphisms
(Lemma \href{https://stacks.math.columbia.edu/tag/01YX}{lemma directed inverse system has limit (Tag 01YX)}).
Let $0 \in I$ and let $f_0 : X_0 \to Y_0$ be a morphism of schemes over $S_0$.
Assume $S_0$, $X_0$, $Y_0$ are quasi-compact and quasi-separated.
Let $f_i : X_i \to Y_i$ be the base change of $f_0$ to $S_i$ and
let $f : X \to Y$ be the base change of $f_0$ to $S$.
\end{situation}

\begin{situation}
\label{situation-local-lefschetz}
Let $(A, \mathfrak m)$ be a Noetherian local ring and $f \in \mathfrak m$.
We set $X = \Spec(A)$ and $X_0 = \Spec(A/fA)$ and we
let $U = X \setminus \{\mathfrak m\}$ and
$U_0 = X_0 \setminus \{\mathfrak m\}$ be the punctured spectrum of
$A$ and $A/fA$.
\end{situation}

\begin{definition}
\label{algebra-definition-standard-etale}
Let $R$ be a ring. Let $g , f  \in R[x]$.
Assume that $f$ is monic and the derivative $f'$ is invertible in
the localization $R[x]_g/(f)$.
In this case the ring map $R \to R[x]_g/(f)$ is said to be
{\it standard \'etale}.
\end{definition}

\begin{lemma}
\label{lemma-reformulate-purity-normal}
Let $(A, \mathfrak m)$ be a Noetherian local ring. Set $X = \Spec(A)$
and let $U = X \setminus \{\mathfrak m\}$. Assume $A$ is normal
of dimension $\geq 2$. The functor
$$
\textit{F\'Et}_U \longrightarrow
\left\{
\begin{matrix}
\text{finite normal }A\text{-algebras }B\text{ such} \\
\text{that }\Spec(B) \to X\text{ is \'etale over }U
\end{matrix}
\right\},
\quad
V \longmapsto \Gamma(V, \mathcal{O}_V)
$$
is an equivalence. Moreover, $V = Y \times_X U$ for some $Y \to X$
finite \'etale if and only if $B = \Gamma(V, \mathcal{O}_V)$
is finite \'etale over $A$.
\end{lemma}

\begin{proof}
Observe that $\text{depth}(A) \geq 2$ because $A$ is normal
(Serre's criterion for normality, Algebra, Lemma
\href{https://stacks.math.columbia.edu/tag/031S}{lemma criterion normal (Tag 031S)}).
Thus the final statement follows from Lemma \href{https://stacks.math.columbia.edu/tag/0BLK}{lemma reformulate purity (Tag 0BLK)}.
Given $\pi : V \to U$ finite \'etale, set $B = \Gamma(V, \mathcal{O}_V)$.
If we can show that $B$ is normal and finite over $A$, then
we obtain the displayed functor. Since there is an obvious
quasi-inverse functor, this is also all that we have to show.

\medskip\noindent
Since $A$ is normal, the scheme $V$ is normal
(Descent, Lemma \href{https://stacks.math.columbia.edu/tag/034F}{lemma normal local smooth (Tag 034F)}).
Hence $V$ is a finite disjoint union of integral schemes
(Properties, Lemma \href{https://stacks.math.columbia.edu/tag/033M}{lemma normal Noetherian (Tag 033M)}).
Thus we may assume $V$ is integral.
In this case the function field $L$ of $V$
(Morphisms, Section \href{https://stacks.math.columbia.edu/tag/01RR}{section rational maps (Tag 01RR)})
is a finite separable extension of the fraction field of $A$
(because we get it by looking at the generic fibre
of $V \to U$ and using Morphisms, Lemma
\href{https://stacks.math.columbia.edu/tag/02GL}{lemma etale over field (Tag 02GL)}).
By Algebra, Lemma
\href{https://stacks.math.columbia.edu/tag/032L}{lemma Noetherian normal domain finite separable extension (Tag 032L)}
the integral closure $B' \subset L$ of $A$ in $L$ is finite over $A$.
By More on Algebra, Lemma \href{https://stacks.math.columbia.edu/tag/0BM4}{lemma integral closure reflexive (Tag 0BM4)}
we see that $B'$ is a reflexive $A$-module, which in turn implies
that $\text{depth}_A(B') \geq 2$ by
More on Algebra, Lemma \href{https://stacks.math.columbia.edu/tag/0AVB}{lemma reflexive over normal (Tag 0AVB)}.

\medskip\noindent
Let $f \in \mathfrak m$. Then $B_f = \Gamma(V \times_U D(f), \mathcal{O}_V)$
(Properties, Lemma \href{https://stacks.math.columbia.edu/tag/01P7}{lemma invert f sections (Tag 01P7)}).
Hence $B'_f = B_f$ because $B_f$ is normal (see above),
finite over $A_f$ with fraction field $L$.
It follows that $V = \Spec(B') \times_X U$.
Then we conclude that $B = B'$ from
Lemma \ref{lemma-sections-over-punctured-spec}
applied to $\Spec(B') \to X$.
This lemma applies because the localizations $B'_{\mathfrak m'}$
of $B'$ at maximal ideals $\mathfrak m' \subset B'$ lying over
$\mathfrak m$ have depth $\geq 2$ by
Algebra, Lemma \href{https://stacks.math.columbia.edu/tag/0AUK}{lemma depth goes down finite (Tag 0AUK)}
and the remark on depth in the preceding paragraph.
\end{proof}


\begin{lemma}
\label{lemma-sections-over-punctured-spec}
Let $(A, \mathfrak m)$ be a Noetherian local ring. Set $X = \Spec(A)$
and let $U = X \setminus \{\mathfrak m\}$.
Let $\pi : Y \to X$ be a finite morphism such that
$\text{depth}(\mathcal{O}_{Y, y}) \geq 2$ for all closed points
$y \in Y$.
Then $Y$ is the spectrum of $B = \mathcal{O}_Y(\pi^{-1}(U))$.
\end{lemma}

\begin{proof}
Set $V = \pi^{-1}(U)$ and denote $\pi' : V \to U$ the restriction of $\pi$.
Consider the $\mathcal{O}_X$-module map
$$
\pi_*\mathcal{O}_Y \longrightarrow j_*\pi'_*\mathcal{O}_V
$$
where $j : U \to X$ is the inclusion morphism. We claim
Divisors, Lemma \href{https://stacks.math.columbia.edu/tag/0E9I}{lemma depth 2 hartog (Tag 0E9I)}
applies to this map. If so, then $B = \Gamma(Y, \mathcal{O}_Y)$
and we see that the lemma holds. Let $x \in X$ be the closed point.
It suffices to show that
$\text{depth}((\pi_*\mathcal{O}_Y)_x) \geq 2$.
Let $y_1, \ldots, y_n \in Y$ be the points mapping to $x$.
By Algebra, Lemma \href{https://stacks.math.columbia.edu/tag/0AUK}{lemma depth goes down finite (Tag 0AUK)}
it suffices to show that
$\text{depth}(\mathcal{O}_{Y, y_i}) \geq 2$ for $i = 1, \ldots, n$.
Since this is the assumption of the lemma the proof is complete.
\end{proof}


\begin{lemma}
\label{lemma-lift-purity}
In Situation \ref{situation-local-lefschetz}.
Let $V$ be finite \'etale over $U$. Assume
\begin{enumerate}
\item $A$ has depth $\geq 3$,
\item $V_0 = V \times_U U_0$ is equal to $Y_0 \times_{X_0} U_0$
for some $Y_0 \to X_0$ finite \'etale.
\end{enumerate}
Then $V = Y \times_X U$ for some $Y \to X$ finite \'etale.
\end{lemma}

\begin{proof}
The assumption of depth forces
$H^1_\mathfrak m(A) = H^2_\mathfrak m(A) = 0$, see
Dualizing Complexes, Lemma \href{https://stacks.math.columbia.edu/tag/0AVZ}{lemma depth (Tag 0AVZ)}.
Hence Lemma \ref{lemma-lift-purity-general} applies.
\end{proof}


\begin{lemma}
\label{lemma-lift-purity-general}
In Situation \ref{situation-local-lefschetz}.
Let $V$ be finite \'etale over $U$. Assume
\begin{enumerate}
\item $H^1_\mathfrak m(A)$ and $H^2_\mathfrak m(A)$
are annihilated by a power of $f$,
\item $V_0 = V \times_U U_0$ is equal to $Y_0 \times_{X_0} U_0$
for some $Y_0 \to X_0$ finite \'etale.
\end{enumerate}
Then $V = Y \times_X U$ for some $Y \to X$ finite \'etale.
\end{lemma}

\begin{proof}
We reduce to the complete case using Lemma \ref{lemma-purity-and-completion}.
(The assumptions carry over; use Dualizing Complexes, Lemma
\href{https://stacks.math.columbia.edu/tag/0ALZ}{lemma torsion change rings (Tag 0ALZ)}.)

\medskip\noindent
In the complete case we can lift $Y_0 \to X_0$ to a finite \'etale
morphism $Y \to X$ by
More on Algebra, Lemma \ref{more-algebra-lemma-finite-etale-equivalence};
observe that $(A, fA)$ is a henselian pair by
More on Algebra, Lemma \href{https://stacks.math.columbia.edu/tag/0ALJ}{lemma complete henselian (Tag 0ALJ)}.
Then we can use Lemma \ref{lemma-fully-faithful-minimal}
to see that $V$ is isomorphic to $Y \times_X U$ and
the proof is complete.
\end{proof}


\begin{lemma}
\label{lemma-purity-and-completion}
Let $(A, \mathfrak m)$ be a Noetherian local ring. Set $X = \Spec(A)$
and let $U = X \setminus \{\mathfrak m\}$.
Let $V$ be finite \'etale over $U$.
Let $A^\wedge$ be the $\mathfrak m$-adic completion of $A$,
let $X' = \Spec(A^\wedge)$ and let $U'$ and $V'$ be the base changes of
$U$ and $V$ to $X'$. The following are equivalent
\begin{enumerate}
\item $V = Y \times_X U$ for some $Y \to X$ finite \'etale, and
\item $V' = Y' \times_{X'} U'$ for some $Y' \to X'$ finite \'etale.
\end{enumerate}
\end{lemma}

\begin{proof}
The implication (1) $\Rightarrow$ (2) follows from taking the base change
of a solution $Y \to X$. Let $Y' \to X'$ be as in (2).
By Lemma \ref{lemma-fill-in-missing} we can find $Y \to X$ finite
such that $V = Y \times_X U$ and $Y' = Y \times_X X'$.
By descent we see that $Y \to X$ is finite \'etale
(Algebra, Lemmas \href{https://stacks.math.columbia.edu/tag/03C4}{lemma descend properties modules (Tag 03C4)} and
\href{https://stacks.math.columbia.edu/tag/00U2}{lemma etale (Tag 00U2)}). This finishes the proof.
\end{proof}


\begin{lemma}
\label{lemma-fill-in-missing}
In Situation \ref{situation-local-lefschetz}. Let $V \to U$ be a finite
morphism.  Let $A^\wedge$ be the $\mathfrak m$-adic completion of $A$,
let $X' = \Spec(A^\wedge)$ and let $U'$ and $V'$ be the base changes of
$U$ and $V$ to $X'$. If $Y' \to X'$ is a finite morphism such that
$V' = Y' \times_{X'} U'$, then there exists a finite morphism $Y \to X$
such that $V = Y \times_X U$ and $Y' = Y \times_X X'$.
\end{lemma}

\begin{proof}
This is a straightforward application of
More on Algebra, Proposition \href{https://stacks.math.columbia.edu/tag/05ER}{proposition equivalence (Tag 05ER)}.
Namely, choose generators $f_1, \ldots, f_t$ of $\mathfrak m$.
For each $i$ write $V \times_U D(f_i) = \Spec(B_i)$.
For $1 \leq i, j \leq n$ we obtain an isomorphism
$\alpha_{ij} : (B_i)_{f_j} \to (B_j)_{f_i}$ of $A_{f_if_j}$-algebras
because the spectrum of both represent $V \times_U D(f_if_j)$.
Write $Y' = \Spec(B')$. Since $V \times_U U' = Y \times_{X'} U'$
we get isomorphisms $\alpha_i : B'_{f_i} \to B_i \otimes_A A^\wedge$.
A straightforward argument shows that $(B', B_i, \alpha_i, \alpha_{ij})$
is an object of $\text{Glue}(A \to A^\wedge, f_1, \ldots, f_t)$, see
More on Algebra, Remark \href{https://stacks.math.columbia.edu/tag/05EL}{remark glueing data (Tag 05EL)}.
Applying the proposition cited above (and using
More on Algebra, Remark \href{https://stacks.math.columbia.edu/tag/05EU}{remark formal glueing algebras (Tag 05EU)}
to obtain the algebra structure) we find an $A$-algebra $B$ such that
$\text{Can}(B)$ is isomorphic to $(B', B_i, \alpha_i, \alpha_{ij})$.
Setting $Y = \Spec(B)$ we see that $Y \to X$ is a morphism
which comes equipped with compatible isomorphisms
$V \cong Y \times_X U$ and $Y' = Y \times_X X'$ as desired.
\end{proof}


\begin{lemma}
\label{lemma-fully-faithful-minimal}
\begin{reference}
\href{https://stacks.math.columbia.edu/bibliography}{Bibliography: Bhatt-local (Corollary 1.11)}
\end{reference}
In Situation \ref{situation-local-lefschetz}. Assume
\begin{enumerate}
\item $H^1_\mathfrak m(A)$ and $H^2_\mathfrak m(A)$ are
annihilated by a power of $f$, and
\item $A$ is henselian or more generally $(A, (f))$ is a henselian pair.
\end{enumerate}
Then the restriction functor
$\textit{F\'Et}_U \longrightarrow \textit{F\'Et}_{U_0}$
is fully faithful.
\end{lemma}

\begin{proof}
By Lemma \href{https://stacks.math.columbia.edu/tag/0BLI}{lemma fully faithful henselian completion (Tag 0BLI)}
we may assume that $A$ is a Noetherian complete local ring.
(The assumptions carry over; use
Dualizing Complexes, Lemma \href{https://stacks.math.columbia.edu/tag/0ALZ}{lemma torsion change rings (Tag 0ALZ)}.)
By Lemma \ref{lemma-restriction-fully-faithful}
the result follows from
Algebraic and Formal Geometry, Lemma
\ref{algebraization-lemma-fully-faithful-alternative}.
\end{proof}


\begin{lemma}
\label{lemma-restriction-fully-faithful}
Let $X$ be a Noetherian scheme and let $Y \subset X$ be a closed subscheme
with ideal sheaf $\mathcal{I} \subset \mathcal{O}_X$.
Assume the completion functor
$$
\textit{Coh}(\mathcal{O}_X)
\longrightarrow 
\textit{Coh}(X, \mathcal{I}),\quad
\mathcal{F} \longmapsto \mathcal{F}^\wedge
$$
is fully faithful on the full subcategory of finite locally free objects
(see above).
Then the restriction functor $\textit{F\'Et}_X \to \textit{F\'Et}_Y$
is fully faithful.
\end{lemma}

\begin{proof}
Since the category of finite \'etale coverings has an
internal hom (Lemma \ref{lemma-internal-hom-finite-etale})
it suffices to prove the following: Given $U$ finite \'etale over $X$
and a morphism $t : Y \to U$ over $X$ there exists a unique section
$s : X \to U$ such that $t = s|_Y$. Picture
$$
\xymatrix{
& U \ar[d]^f \\
Y \ar[r] \ar[ru] & X \ar@{..>}@/^1em/[u]
}
$$
Finding the dotted arrow $s$ is the same thing as finding an
$\mathcal{O}_X$-algebra map
$$
s^\sharp : f_*\mathcal{O}_U \longrightarrow \mathcal{O}_X
$$
which reduces modulo the ideal sheaf of $Y$ to the given algebra map
$t^\sharp : f_*\mathcal{O}_U \to \mathcal{O}_Y$.
By Lemma \ref{lemma-thickening} we can lift $t$ uniquely to a compatible
system of maps $t_n : Y_n \to U$ and hence a map
$$
\lim t_n^\sharp : f_*\mathcal{O}_U \longrightarrow \lim \mathcal{O}_{Y_n}
$$
of sheaves of algebras on $X$.
Since $f_*\mathcal{O}_U$ is a finite locally free $\mathcal{O}_X$-module,
we conclude that we get a unique $\mathcal{O}_X$-module map
$\sigma : f_*\mathcal{O}_U \to \mathcal{O}_X$ whose completion
is $\lim t_n^\sharp$. To see that $\sigma$ is an algebra homomorphism,
we need to check that the diagram
$$
\xymatrix{
f_*\mathcal{O}_U \otimes_{\mathcal{O}_X} f_*\mathcal{O}_U
\ar[r] \ar[d]_{\sigma \otimes \sigma} &
f_*\mathcal{O}_U \ar[d]^\sigma \\
\mathcal{O}_X \otimes_{\mathcal{O}_X} \mathcal{O}_X \ar[r] &
\mathcal{O}_X
}
$$
commutes. For every $n$ we know this diagram commutes after restricting
to $Y_n$, i.e., the diagram commutes after applying the completion functor.
Hence by faithfulness of the completion functor we conclude.
\end{proof}


\begin{lemma}
\label{algebraization-lemma-fully-faithful-alternative}
Let $A$ be a Noetherian ring. Let $f \in \mathfrak a \subset A$
be an element of an ideal of $A$. Let $U = \Spec(A) \setminus V(\mathfrak a)$.
Assume
\begin{enumerate}
\item $A$ is $f$-adically complete,
\item $H^1_\mathfrak a(A)$ and $H^2_\mathfrak a(A)$ are
annihilated by a power of $f$.
\end{enumerate}
Then the completion functor
$$
\textit{Coh}(\mathcal{O}_U)
\longrightarrow
\textit{Coh}(U, I\mathcal{O}_U),
\quad
\mathcal{F} \longmapsto \mathcal{F}^\wedge
$$
is fully faithful on the full subcategory of
finite locally free objects.
\end{lemma}

\begin{proof}
By Lemma \ref{algebraization-lemma-completion-fully-faithful}
it suffices to show that
$$
\Gamma(U, \mathcal{O}_U) =
\lim \Gamma(U, \mathcal{O}_U/I^n\mathcal{O}_U)
$$
This follows immediately from
Lemma \ref{algebraization-lemma-alternative-H0}.
\end{proof}


\begin{lemma}
\label{algebraization-lemma-alternative-H0}
Let $A$ be a Noetherian ring. Let $f \in \mathfrak a \subset A$
be an element of an ideal of $A$. Let $M$ be a finite $A$-module.
Assume
\begin{enumerate}
\item $A$ is $f$-adically complete,
\item $H^1_\mathfrak a(M)$ and $H^2_\mathfrak a(M)$ are
annihilated by a power of $f$.
\end{enumerate}
Then with $U = \Spec(A) \setminus V(\mathfrak a)$ the map
$$
\Gamma(U, \widetilde{M})
\longrightarrow
\lim \Gamma(U, \widetilde{M/f^nM})
$$
is an isomorphism.
\end{lemma}

\begin{proof}
We may apply Lemma \href{https://stacks.math.columbia.edu/tag/0BLD}{lemma formal functions principal (Tag 0BLD)} to $U$ and
$\mathcal{F} = \widetilde{M}|_U$ because $\mathcal{F}$ is a Noetherian object in
the category of coherent $\mathcal{O}_U$-modules.
Since $H^1(U, \mathcal{F}) = H^2_\mathfrak a(M)$
(Local Cohomology, Lemma
\href{https://stacks.math.columbia.edu/tag/0BK0}{lemma finiteness pushforwards and H1 local (Tag 0BK0)})
is annihilated by a power of $f$, we see that
its $f$-adic Tate module is zero.
Hence the lemma shows $\lim H^0(U, \mathcal{F}/f^n \mathcal{F})$
is equal to the usual $f$-adic completion of $H^0(U, \mathcal{F})$.
Consider the short exact sequence
$$
0 \to M/H^0_\mathfrak a(M) \to H^0(U, \mathcal{F}) \to H^1_\mathfrak a(M) \to 0
$$
of Local Cohomology, Lemma
\href{https://stacks.math.columbia.edu/tag/0BK0}{lemma finiteness pushforwards and H1 local (Tag 0BK0)}.
Since $M/H^0_\mathfrak a(M)$ is a finite $A$-module, it is
complete, see Algebra, Lemma \href{https://stacks.math.columbia.edu/tag/00MA}{lemma completion tensor (Tag 00MA)}.
Since $H^1_\mathfrak a(M)$ is killed by a power of $f$, we conclude
from Algebra, Lemma \href{https://stacks.math.columbia.edu/tag/0BNG}{lemma completion differ by torsion (Tag 0BNG)}
that $H^0(U, \mathcal{F})$ is complete as well. This finishes the proof.
\end{proof}


\begin{lemma}
\label{algebraization-lemma-completion-fully-faithful}
Let $X$ be a Noetherian scheme and let $Y \subset X$ be a closed subscheme.
Let $Y_n \subset X$ be the $n$th infinitesimal neighbourhood of $Y$ in $X$.
Consider the following conditions
\begin{enumerate}
\item $X$ is quasi-affine and
$\Gamma(X, \mathcal{O}_X) \to \lim \Gamma(Y_n, \mathcal{O}_{Y_n})$
is an isomorphism,
\item $X$ has an ample invertible module $\mathcal{L}$ and
$\Gamma(X, \mathcal{L}^{\otimes m}) \to
\lim \Gamma(Y_n, \mathcal{L}^{\otimes m}|_{Y_n})$
is an isomorphism for all $m \gg 0$,
\item for every finite locally free $\mathcal{O}_X$-module
$\mathcal{E}$ the map
$\Gamma(X, \mathcal{E}) \to \lim \Gamma(Y_n, \mathcal{E}|_{Y_n})$
is an isomorphism, and
\item the completion functor
$\textit{Coh}(\mathcal{O}_X) \to \textit{Coh}(X, \mathcal{I})$
is fully faithful on the full subcategory of finite locally free
objects.
\end{enumerate}
Then (1) $\Rightarrow$ (2) $\Rightarrow$ (3) $\Rightarrow$ (4)
and (4) $\Rightarrow$ (3).
\end{lemma}

\begin{proof}
Proof of (3) $\Rightarrow$ (4). If $\mathcal{F}$ and $\mathcal{G}$
are finite locally free on $X$, then considering
$\mathcal{H} = \SheafHom_{\mathcal{O}_X}(\mathcal{G}, \mathcal{F})$
and using Cohomology of Schemes, Lemma
\href{https://stacks.math.columbia.edu/tag/0882}{lemma completion internal hom (Tag 0882)}
we see that (3) implies (4).

\medskip\noindent
Proof of (2) $\rightarrow$ (3). Namely, let $\mathcal{L}$ be ample
on $X$ and suppose that $\mathcal{E}$ is a
finite locally free $\mathcal{O}_X$-module.
We claim we can find a universally exact sequence
$$
0 \to \mathcal{E} \to
(\mathcal{L}^{\otimes p})^{\oplus r} \to
(\mathcal{L}^{\otimes q})^{\oplus s}
$$
for some $r, s \geq 0$ and $0 \ll p \ll q$. If this holds, then
using the exact sequence
$$
0 \to \lim \Gamma(\mathcal{E}|_{Y_n}) \to
\lim \Gamma((\mathcal{L}^{\otimes p})^{\oplus r}|_{Y_n}) \to
\lim \Gamma((\mathcal{L}^{\otimes q})^{\oplus s}|_{Y_n})
$$
and the isomorphisms in (2) we get the isomorphism in (3).
To prove the claim, consider the dual locally free module
$\SheafHom_{\mathcal{O}_X}(\mathcal{E}, \mathcal{O}_X)$
and apply
Properties, Proposition \href{https://stacks.math.columbia.edu/tag/01Q3}{proposition characterize ample (Tag 01Q3)}
to find a surjection
$$
(\mathcal{L}^{\otimes -p})^{\oplus r}
\longrightarrow
\SheafHom_{\mathcal{O}_X}(\mathcal{E}, \mathcal{O}_X)
$$
Taking duals we obtain the first map in the exact sequence
(it is universally injective because being a surjection is universal).
Repeat with the cokernel to get the second. Some details omitted.

\medskip\noindent
Proof of (1) $\Rightarrow$ (2). This is true because if $X$ is quasi-affine
then $\mathcal{O}_X$ is an ample invertible module, see
Properties, Lemma \href{https://stacks.math.columbia.edu/tag/01QE}{lemma quasi affine O ample (Tag 01QE)}.

\medskip\noindent
We omit the proof of (4) $\Rightarrow$ (3).
\end{proof}


\begin{lemma}
\label{lemma-ramification-quasi-finite-flat}
Let $f : X \to Y$ be a morphism of locally Noetherian schemes.
Let $x \in X$. Assume
\begin{enumerate}
\item $f$ is flat,
\item $f$ is quasi-finite at $x$,
\item $x$ is not a generic point of an irreducible component of $X$,
\item for specializations $x' \leadsto x$ with
$\dim(\mathcal{O}_{X, x'}) = 1$ our $f$ is unramified at $x'$.
\end{enumerate}
Then $f$ is \'etale at $x$.
\end{lemma}

\begin{proof}
Observe that the set of points where $f$ is unramified is the same as
the set of points where $f$ is \'etale and that this set is open.
See Morphisms, Definitions \ref{morphisms-definition-unramified}
and \ref{morphisms-definition-etale} and
Lemma \href{https://stacks.math.columbia.edu/tag/02GV}{lemma flat unramified etale (Tag 02GV)}.
To check $f$ is \'etale at $x$ we may work \'etale
locally on the base and on the
target (Descent, Lemmas \href{https://stacks.math.columbia.edu/tag/02VN}{lemma descending property etale (Tag 02VN)} and
\href{https://stacks.math.columbia.edu/tag/036W}{lemma etale etale local source (Tag 036W)}).
Thus we can apply More on Morphisms, Lemma
\href{https://stacks.math.columbia.edu/tag/02LK}{lemma etale makes quasi finite finite at point (Tag 02LK)}
and assume that $f : X \to Y$ is finite and that $x$ is the unique
point of $X$ lying over $y = f(x)$.
Then it follows that $f$ is finite locally free
(Morphisms, Lemma \href{https://stacks.math.columbia.edu/tag/02KB}{lemma finite flat (Tag 02KB)}).

\medskip\noindent
Assume $f$ is finite locally free and that $x$ is the unique point of
$X$ lying over $y = f(x)$. By
Discriminants, Lemma \href{https://stacks.math.columbia.edu/tag/0BJF}{lemma discriminant (Tag 0BJF)}
we find a locally principal closed subscheme $D_\pi \subset Y$
such that $y' \in D_\pi$ if and only if there exists an $x' \in X$
with $f(x') = y'$ and $f$ ramified at $x'$. Thus we have to prove
that $y \not \in D_\pi$. Assume $y \in D_\pi$ to get a contradiction.

\medskip\noindent
By condition (3) we have $\dim(\mathcal{O}_{X, x}) \geq 1$.
We have $\dim(\mathcal{O}_{X, x}) = \dim(\mathcal{O}_{Y, y})$ by
Algebra, Lemma \href{https://stacks.math.columbia.edu/tag/00ON}{lemma dimension base fibre equals total (Tag 00ON)}.
By Lemma \href{https://stacks.math.columbia.edu/tag/0BJG}{lemma find point codim 1 (Tag 0BJG)}
we can find $y' \in D_\pi$ specializing to $y$
with $\dim(\mathcal{O}_{Y, y'}) = 1$.
Choose $x' \in X$ with $f(x') = y'$ where $f$ is ramified. Since $f$
is finite it is closed, and hence $x' \leadsto x$.
We have $\dim(\mathcal{O}_{X, x'}) = \dim(\mathcal{O}_{Y, y'}) = 1$
as before. This contradicts property (4).
\end{proof}


\begin{lemma}
\label{lemma-internal-hom-finite-etale}
Let $X$ be a scheme. Given $U, V$ finite \'etale over $X$ there
exists a scheme $W$ finite \'etale over $X$ such that
$$
\Mor_X(X, W) = \Mor_X(U, V)
$$
and such that the same remains true after any base change.
\end{lemma}

\begin{proof}
By More on Morphisms, Lemma
\href{https://stacks.math.columbia.edu/tag/0BL5}{lemma hom from finite locally free separated lqf (Tag 0BL5)}
there exists a scheme $W$ representing $\mathit{Mor}_X(U, V)$.
(Use that an \'etale morphism is locally quasi-finite by
Morphisms, Lemmas \href{https://stacks.math.columbia.edu/tag/03WS}{lemma etale locally quasi finite (Tag 03WS)}
and that a finite morphism is separated.)
This scheme clearly satisfies the formula after any base change.
To finish the proof we have to show that $W \to X$ is finite \'etale.
This we may do after replacing $X$ by the members of an \'etale
covering (Descent, Lemmas \href{https://stacks.math.columbia.edu/tag/02LA}{lemma descending property finite (Tag 02LA)}
and \href{https://stacks.math.columbia.edu/tag/02KU}{lemma descending property separated (Tag 02KU)}).
Thus by \'Etale Morphisms, Lemma \href{https://stacks.math.columbia.edu/tag/04HN}{lemma finite etale etale local (Tag 04HN)}
we may assume that $U = \coprod_{i = 1, \ldots, n} X$
and $V = \coprod_{j = 1, \ldots, m} X$.
In this case
$W = \coprod_{\alpha : \{1, \ldots, n\} \to \{1, \ldots, m\}} X$
by inspection (details omitted) and the proof is complete.
\end{proof}


\begin{lemma}
\label{lemma-thickening}
Let $X \subset X'$ be a thickening of schemes. The functor
$$
\textit{F\'Et}_{X'} \longrightarrow \textit{F\'Et}_X,\quad
U' \longmapsto U' \times_{X'} X
$$
is an equivalence of categories.
\end{lemma}

\begin{proof}
For a discussion of thickenings see
More on Morphisms, Section \href{https://stacks.math.columbia.edu/tag/04EW}{section thickenings (Tag 04EW)}.
Let $U' \to X'$ be an \'etale morphism such that $U = U' \times_{X'} X \to X$
is finite \'etale. Then $U' \to X'$ is finite \'etale as well.
This follows for example from More on Morphisms, Lemma
\href{https://stacks.math.columbia.edu/tag/0BPG}{lemma properties that extend over thickenings (Tag 0BPG)}.
Now, if $X \subset X'$ is a finite order thickening then this remark
combined with \'Etale Morphisms, Theorem
\href{https://stacks.math.columbia.edu/tag/039R}{theorem remarkable equivalence (Tag 039R)}
proves the lemma. Below we will prove the lemma for general thickenings, but
we suggest the reader skip the proof.

\medskip\noindent
Let $X' = \bigcup X_i'$ be an affine open covering. Set
$X_i = X \times_{X'} X_i'$, $X_{ij}' = X'_i \cap X'_j$,
$X_{ij} = X \times_{X'} X_{ij}'$, $X_{ijk}' = X'_i \cap X'_j \cap X'_k$,
$X_{ijk} = X \times_{X'} X_{ijk}'$.
Suppose that we can prove
the theorem for each of the thickenings
$X_i \subset X'_i$, $X_{ij} \subset X_{ij}'$, and $X_{ijk} \subset X_{ijk}'$.
Then the result follows for $X \subset X'$ by relative glueing of
schemes, see
Constructions, Section \href{https://stacks.math.columbia.edu/tag/01LG}{section relative glueing (Tag 01LG)}.
Observe that the schemes $X_i'$, $X_{ij}'$, $X_{ijk}'$ are
each separated as open subschemes of affine schemes. Repeating the
argument one more time we reduce to the case where the schemes
$X'_i$, $X_{ij}'$, $X_{ijk}'$ are affine.

\medskip\noindent
In the affine case we have $X' = \Spec(A')$ and $X = \Spec(A'/I')$
where $I'$ is a locally nilpotent ideal. Then $(A', I')$ is a
henselian pair (More on Algebra, Lemma
\href{https://stacks.math.columbia.edu/tag/0ALI}{lemma locally nilpotent henselian (Tag 0ALI)})
and the result follows from Lemma \href{https://stacks.math.columbia.edu/tag/09ZS}{lemma gabber (Tag 09ZS)} (which is
much easier in this case).
\end{proof}


\begin{lemma}
\label{more-algebra-lemma-finite-etale-equivalence}
Let $(A, I)$ be a henselian pair. The functor $B \to B/IB$ determines
an equivalence between finite \'etale $A$-algebras and finite \'etale
$A/I$-algebras.
\end{lemma}

\begin{proof}
Let $B, B'$ be two $A$-algebras finite \'etale over $A$.
Then $B' \to B'' = B \otimes_A B'$ is finite \'etale as well
(Algebra, Lemmas \href{https://stacks.math.columbia.edu/tag/00U2}{lemma etale (Tag 00U2)} and
\href{https://stacks.math.columbia.edu/tag/02JK}{lemma base change integral (Tag 02JK)}).
Now we have $1$-to-$1$ correspondences between
\begin{enumerate}
\item $A$-algebra maps $B \to B'$,
\item sections of $B' \to B''$, and
\item idempotents $e$ of $B''$ such that $B' \to B'' \to eB''$ is
an isomorphism.
\end{enumerate}
The bijection between (2) and (3) sends $\sigma : B'' \to B'$
to $e$ such that $(1 - e)$ is the idempotent
that generates the kernel of $\sigma$ which exists by
Algebra, Lemmas \href{https://stacks.math.columbia.edu/tag/00U7}{lemma map between etale (Tag 00U7)} and
\href{https://stacks.math.columbia.edu/tag/00U8}{lemma surjective flat finitely presented (Tag 00U8)}.
There is a similar correspondence between
$A/I$-algebra maps $B/IB \to B'/IB'$ and idempotents
$\overline{e}$ of $B''/IB''$ such that
$B'/IB' \to B''/IB'' \to \overline{e}(B''/IB'')$ is
an isomorphism. However every idempotent $\overline{e}$ of $B''/IB''$
lifts uniquely to an idempotent $e$ of $B''$
(Lemma \href{https://stacks.math.columbia.edu/tag/09XI}{lemma characterize henselian pair (Tag 09XI)}).
Moreover, if $B'/IB' \to \overline{e}(B''/IB'')$ is an isomorphism,
then $B' \to eB''$ is an isomorphism too by Nakayama's lemma
(Algebra, Lemma \href{https://stacks.math.columbia.edu/tag/00DV}{lemma NAK (Tag 00DV)}).
In this way we see that the functor is fully faithful.

\medskip\noindent
Essential surjectivity. Let $A/I \to C$ be a finite \'etale map.
By Algebra, Lemma \href{https://stacks.math.columbia.edu/tag/04D1}{lemma lift etale (Tag 04D1)}
there exists an \'etale map $A \to B$ such that $B/IB \cong C$.
Let $B'$ be the integral closure of $A$ in $B$. 
By Lemma \href{https://stacks.math.columbia.edu/tag/09XH}{lemma helper finite type (Tag 09XH)} we have
$B'/IB' = C \times C'$ for some ring $C'$
and $B'_g \cong B_g$ for some $g \in B'$ mapping to $(1, 0) \in C \times C'$.
Since idempotents lift
(Lemma \href{https://stacks.math.columbia.edu/tag/09XI}{lemma characterize henselian pair (Tag 09XI)})
we get $B' = B'_1 \times B'_2$ with $C = B'_1/IB'_1$ and $C' = B'_2/IB'_2$.
The image of $g$ in $B'_1$ is invertible.
Then $B_g = B'_g = B'_1 \times (B_2)_g$ and this implies
that $A \to B'_1$ is \'etale.
We conclude that $B'_1$ is finite \'etale over $A$
(integral \'etale implies finite \'etale by
Algebra, Lemma \href{https://stacks.math.columbia.edu/tag/02JJ}{lemma characterize finite in terms of integral (Tag 02JJ)}
for example)
and the proof is done.
\end{proof}


\begin{lemma}
\label{lemma-local-purity}
Let $(A, \mathfrak m)$ be a regular local ring of dimension $d \geq 2$.
Set $X = \Spec(A)$ and $U = X \setminus \{\mathfrak m\}$. Then
\begin{enumerate}
\item the functor $\textit{F\'Et}_X \to \textit{F\'Et}_U$
is essentially surjective, i.e., purity holds for $A$,
\item any finite $A \to B$ with $B$ normal which
induces a finite \'etale morphism on punctured spectra is \'etale.
\end{enumerate}
\end{lemma}

\begin{proof}
Recall that a regular local ring is normal by
Algebra, Lemma \href{https://stacks.math.columbia.edu/tag/0567}{lemma regular normal (Tag 0567)}.
Hence (1) and (2) are equivalent by
Lemma \ref{lemma-reformulate-purity-normal}.
We prove the lemma by induction on $d$.

\medskip\noindent
The case $d = 2$. In this case $A \to B$ is flat.
Namely, we have going down for $A \to B$ by
Algebra, Proposition \href{https://stacks.math.columbia.edu/tag/00H8}{proposition going down normal integral (Tag 00H8)}.
Then $\dim(B_{\mathfrak m'}) = 2$ for all maximal ideals
$\mathfrak m' \subset B$ by
Algebra, Lemma \href{https://stacks.math.columbia.edu/tag/00ON}{lemma dimension base fibre equals total (Tag 00ON)}.
Then $B_{\mathfrak m'}$ is Cohen-Macaulay by
Algebra, Lemma \href{https://stacks.math.columbia.edu/tag/031S}{lemma criterion normal (Tag 031S)}.
Hence and this is the important step
Algebra, Lemma \href{https://stacks.math.columbia.edu/tag/00R4}{lemma CM over regular flat (Tag 00R4)}
applies to show $A \to B_{\mathfrak m'}$ is flat.
Then Algebra, Lemma \href{https://stacks.math.columbia.edu/tag/00HT}{lemma flat localization (Tag 00HT)}
shows $A \to B$ is flat. Thus we can apply
Lemma \ref{lemma-ramification-quasi-finite-flat}
(or you can directly argue using the easier
Discriminants, Lemma \href{https://stacks.math.columbia.edu/tag/0BJF}{lemma discriminant (Tag 0BJF)})
to see that $A \to B$ is \'etale.

\medskip\noindent
The case $d \geq 3$. Let $V \to U$ be finite \'etale.
Let $f \in \mathfrak m_A$, $f \not \in \mathfrak m_A^2$.
Then $A/fA$ is a regular local ring of dimension $d - 1 \geq 2$, see
Algebra, Lemma \href{https://stacks.math.columbia.edu/tag/00NQ}{lemma regular ring CM (Tag 00NQ)}.
Let $U_0$ be the punctured spectrum of $A/fA$ and let
$V_0 = V \times_U U_0$.
By Lemma \ref{lemma-lift-purity}
it suffices to show that $V_0$ is in the essential
image of $\textit{F\'Et}_{\Spec(A/fA)} \to \textit{F\'Et}_{U_0}$.
This follows from the induction hypothesis.
\end{proof}


\begin{lemma}[Purity of branch locus]
\label{lemma-purity}
\begin{reference}
\href{https://stacks.math.columbia.edu/bibliography}{Bibliography: Nagata-Purity} and \href{https://stacks.math.columbia.edu/bibliography}{Bibliography: SGA1 (Exp. X, Thm. 3.1)}
\end{reference}
\begin{history}
This result was first stated and proved by Zariski in
geometric form in \href{https://stacks.math.columbia.edu/bibliography}{Bibliography: Zariski-Purity}.
The generalization to nonperfect ground fields by Nagata
was published as the next article in the same volume of the
Proceedings of the National Academy of Sciences of the United States of America
in \href{https://stacks.math.columbia.edu/bibliography}{Bibliography: Nagata-Remarks-Purity}. In the following year Nagata
proved the result for Noetherian local rings in \href{https://stacks.math.columbia.edu/bibliography}{Bibliography: Nagata-Purity}.
His proof uses a result of Chow which is a Bertini theorem for
complete local domains, see \href{https://stacks.math.columbia.edu/bibliography}{Bibliography: Chow-Bertini};
the history of Bertini's theorems is discussed in
Kleiman's historical article \href{https://stacks.math.columbia.edu/bibliography}{Bibliography: Kleiman-Bertini}.
A few years later a completely different proof was found by
Auslander, see \href{https://stacks.math.columbia.edu/bibliography}{Bibliography: Auslander-Purity}.
\end{history}
Let $f : X \to Y$ be a morphism of locally Noetherian schemes.
Let $x \in X$ and set $y = f(x)$. Assume
\begin{enumerate}
\item $\mathcal{O}_{X, x}$ is normal,
\item $\mathcal{O}_{Y, y}$ is regular,
\item $f$ is quasi-finite at $x$,
\item $\dim(\mathcal{O}_{X, x}) = \dim(\mathcal{O}_{Y, y}) \geq 1$
\item for specializations $x' \leadsto x$ with
$\dim(\mathcal{O}_{X, x'}) = 1$ our $f$ is unramified at $x'$.
\end{enumerate}
Then $f$ is \'etale at $x$.
\end{lemma}

\begin{proof}
We will prove the lemma by induction on
$d = \dim(\mathcal{O}_{X, x}) = \dim(\mathcal{O}_{Y, y})$.

\medskip\noindent
An uninteresting case is when $d = 1$.
In that case we are assuming that $f$ is unramified at $x$
and that $\mathcal{O}_{Y, y}$ is a discrete valuation ring
(Algebra, Lemma \href{https://stacks.math.columbia.edu/tag/00PD}{lemma characterize dvr (Tag 00PD)}).
Then $\mathcal{O}_{X, x}$ is flat over $\mathcal{O}_{Y, y}$
(otherwise the map would not be quasi-finite at $x$)
and we see that $f$ is flat at $x$. Since flat $+$
unramified is \'etale we conclude (some details omitted).

\medskip\noindent
The case $d \geq 2$. We will use induction on $d$ to reduce
to the case discussed in Lemma \ref{lemma-local-purity}.
To check $f$ is \'etale at $x$ we may work \'etale locally
on the base and on the target
(Descent, Lemmas \href{https://stacks.math.columbia.edu/tag/02VN}{lemma descending property etale (Tag 02VN)} and
\href{https://stacks.math.columbia.edu/tag/036W}{lemma etale etale local source (Tag 036W)}).
Thus we can apply More on Morphisms, Lemma
\href{https://stacks.math.columbia.edu/tag/02LK}{lemma etale makes quasi finite finite at point (Tag 02LK)}
and assume that $f : X \to Y$ is finite and that $x$ is the unique
point of $X$ lying over $y$. Here we use that \'etale extensions of
local rings do not change dimension, normality, and regularity, see
More on Algebra, Section \href{https://stacks.math.columbia.edu/tag/0AGY}{section permanence etale (Tag 0AGY)}
and
\'Etale Morphisms, Section \href{https://stacks.math.columbia.edu/tag/025L}{section properties permanence (Tag 025L)}.

\medskip\noindent
Next, we can base change by $\Spec(\mathcal{O}_{Y, y})$
and assume that $Y$ is the spectrum of a regular local ring.
It follows that $X = \Spec(\mathcal{O}_{X, x})$ as
every point of $X$ necessarily specializes to $x$.

\medskip\noindent
The ring map $\mathcal{O}_{Y, y} \to \mathcal{O}_{X, x}$ is
finite and necessarily injective (by equality of dimensions).
We conclude we have going down for
$\mathcal{O}_{Y, y} \to \mathcal{O}_{X, x}$ by
Algebra, Proposition \href{https://stacks.math.columbia.edu/tag/00H8}{proposition going down normal integral (Tag 00H8)}
(and the fact that a regular ring is a normal ring by
Algebra, Lemma \href{https://stacks.math.columbia.edu/tag/0567}{lemma regular normal (Tag 0567)}).
Pick $x' \in X$, $x' \not = x$ with image $y' = f(x')$.
Then $\mathcal{O}_{X, x'}$ is normal as a localization
of a normal domain. Similarly, $\mathcal{O}_{Y, y'}$ is
regular (see Algebra, Lemma
\href{https://stacks.math.columbia.edu/tag/0AFS}{lemma localization of regular local is regular (Tag 0AFS)}).
We have $\dim(\mathcal{O}_{X, x'}) = \dim(\mathcal{O}_{Y, y'})$ by
Algebra, Lemma \href{https://stacks.math.columbia.edu/tag/00ON}{lemma dimension base fibre equals total (Tag 00ON)}
(we checked going down above).
Of course these dimensions are strictly less than $d$ as $x' \not = x$
and by induction on $d$ we conclude that $f$ is \'etale at $x'$.

\medskip\noindent
Thus we arrive at the following situation: We have a finite
local homomorphism $A \to B$ of Noetherian local rings
of dimension $d \geq 2$, with $A$ regular, $B$ normal, which
induces a finite \'etale morphism $V \to U$ on punctured spectra.
Our goal is to show that $A \to B$ is \'etale.
This follows from Lemma \ref{lemma-local-purity}
and the proof is complete.
\end{proof}


\begin{lemma}
\label{lemma-key-purity-ramification}
Let $f : X \to \Spec(A)$ be a finite type morphism.
Let $x \in X$ be a point. Assume
\begin{enumerate}
\item $A$ is an excellent regular local ring,
\item $\mathcal{O}_{X, x}$ is normal of dimension $2$,
\item $f$ is \'etale outside of $\overline{\{x\}}$.
\end{enumerate}
Then $f$ is \'etale at $x$.
\end{lemma}

\begin{proof}
We first replace $X$ by an affine open neighbourhood of $x$.
Observe that $\mathcal{O}_{X, x}$ is an excellent local ring
(More on Algebra, Lemma \href{https://stacks.math.columbia.edu/tag/07QU}{lemma finite type over excellent (Tag 07QU)}).
Thus we can choose a minimal resolution of singularities
$W \to \Spec(\mathcal{O}_{X, x})$, see
Resolution of Surfaces, Theorem \href{https://stacks.math.columbia.edu/tag/0BGP}{theorem resolve (Tag 0BGP)}.
After possibly replacing $X$ by an affine open neighbourhood of $x$
we can find a proper morphism $b : X' \to X$ such that
$X' \times_X \Spec(\mathcal{O}_{X, x}) = W$, see
Limits, Lemma \ref{limits-lemma-glueing-near-closed-point}.
After shrinking $X$ further, we may assume $X'$ is regular.
Namely, we know $W$ is regular and $X'$ is excellent
and the regular locus of the spectrum of an excellent ring is open.
Since $W \to \Spec(\mathcal{O}_{X, x})$ is projective
(as a sequence of normalized blowing ups), we may assume
after shrinking $X$ that $b$ is projective (details omitted).
Let $U = X \setminus \overline{\{x\}}$.
Since $W \to \Spec(\mathcal{O}_{X, x})$ is an isomorphism over the
punctured spectrum, we may assume $b : X' \to X$ is an isomorphism over $U$.
Thus we may and will think of $U$ as an open subscheme of $X'$ as well.
Set $f' = f \circ b : X' \to \Spec(A)$.

\medskip\noindent
Since $A$ is regular we see that $\mathcal{O}_Y$ is a dualizing complex for $Y$.
Hence $f^!\mathcal{O}_Y$ is a dualzing complex on $X$
(Duality for Schemes, Lemma \href{https://stacks.math.columbia.edu/tag/0AA3}{lemma shriek dualizing (Tag 0AA3)}).
The Cohen-Macaulay locus of $X$ is open by
Duality for Schemes, Lemma \href{https://stacks.math.columbia.edu/tag/0AWT}{lemma dualizing module CM scheme (Tag 0AWT)}
(this can also be proven using excellency).
Since $\mathcal{O}_{X, x}$ is Cohen-Macaulay, after shrinking
$X$ we may assume $X$ is Cohen-Macaulay.
Observe that an \'etale morphism is a local complete intersection.
Thus
Duality for Schemes, Lemma \ref{duality-lemma-fundamental-class-almost-lci}
applies with $r = 0$ and we get a map
$$
\mathcal{O}_X \longrightarrow \omega_{X/Y} = H^0(f^!\mathcal{O}_Y)
$$
which is an isomorphism over $X \setminus \overline{\{x\}}$.
Since $\omega_{X/Y}$ is $(S_2)$ by
Duality for Schemes, Lemma \href{https://stacks.math.columbia.edu/tag/0E9V}{lemma shriek over CM (Tag 0E9V)}
we find this map is an isomorphism by
Divisors, Lemma \href{https://stacks.math.columbia.edu/tag/0AVM}{lemma check isomorphism via depth and ass (Tag 0AVM)}.
This already shows that $X$ and in particular $\mathcal{O}_{X, x}$ is
Gorenstein.

\medskip\noindent
Set $\omega_{X'/Y} = H^0((f')^!\mathcal{O}_Y)$. Arguing in exactly the
same manner as above we find that $(f')^!\mathcal{O}_Y = \omega_{X'/Y}[0]$
is a dualizing complex for $X'$. Since $X'$ is regular
the morphism $X' \to Y$ is a local complete intersection morphism, see
More on Morphisms, Lemma
\href{https://stacks.math.columbia.edu/tag/0E9K}{lemma morphism regular schemes is lci (Tag 0E9K)}.
By Duality for Schemes, Lemma \ref{duality-lemma-fundamental-class-lci}
there exists a map
$$
\mathcal{O}_{X'} \longrightarrow \omega_{X'/Y}
$$
which is an isomorphism over $U$. We conclude
$\omega_{X'/Y} = \mathcal{O}_{X'}(E)$ for some effective
Cartier divisor $E \subset X'$ disjoint from $U$.

\medskip\noindent
Since $\omega_{X/Y} = \mathcal{O}_Y$ we see that
$\omega_{X'/Y} = b^! f^!\mathcal{O}_Y = b^!\mathcal{O}_X$.
Returning to $W \to \Spec(\mathcal{O}_{X, x})$
we see that $\omega_W = \mathcal{O}_W(E|_W)$.
By Lemma \ref{lemma-characterize-rational-singularity}
we find $E|_W = 0$.
This means that $f' : X' \to Y$ is \'etale by (the already used)
Duality for Schemes, Lemma \ref{duality-lemma-fundamental-class-lci}.
This immediately finishes the proof, as \'etaleness
of $f'$ forces $b$ to be an isomorphism.
\end{proof}


\begin{lemma}
\label{lemma-characterize-rational-singularity}
Let $A$ be a Noetherian normal local domain of dimension $2$.
Assume $A$ is Nagata, has a dualizing module $\omega_A$, and has a
resolution of singularities $f : X \to \Spec(A)$.
Let $\omega_X$ be as in Resolution of Surfaces,
Remark \href{https://stacks.math.columbia.edu/tag/0B4R}{remark dualizing setup (Tag 0B4R)}.
If $\omega_X \cong \mathcal{O}_X(E)$ for some effective
Cartier divisor $E \subset X$ supported on the exceptional
fibre, then $A$ defines a rational singularity.
If $f$ is a minimal resolution, then $E = 0$.
\end{lemma}

\begin{proof}
There is a trace map $Rf_*\omega_X \to \omega_A$, see
Duality for Schemes, Section \href{https://stacks.math.columbia.edu/tag/0AWG}{section trace (Tag 0AWG)}.
By Grauert-Riemenschneider
(Resolution of Surfaces,
Proposition \href{https://stacks.math.columbia.edu/tag/0AXD}{proposition Grauert Riemenschneider (Tag 0AXD)})
we have $R^1f_*\omega_X = 0$.
Thus the trace map is a map $f_*\omega_X \to \omega_A$.
Then we can consider
$$
\mathcal{O}_{\Spec(A)} = f_*\mathcal{O}_X \to f_*\omega_X \to \omega_A
$$
where the first map comes from the map
$\mathcal{O}_X \to \mathcal{O}_X(E) = \omega_X$ which is
assumed to exist in the statement of the lemma.
The composition is an isomorphism by Divisors, Lemma
\href{https://stacks.math.columbia.edu/tag/0AVM}{lemma check isomorphism via depth and ass (Tag 0AVM)}
as it is an isomorphism over the punctured spectrum of $A$
(by the assumption in the lemma and the fact that $f$ is an isomorphism
over the punctured spectrum) and $A$ and $\omega_A$
are $A$-modules of depth $2$ (by
Algebra, Lemma \href{https://stacks.math.columbia.edu/tag/031S}{lemma criterion normal (Tag 031S)} and
Dualizing Complexes, Lemma \href{https://stacks.math.columbia.edu/tag/0AWE}{lemma depth dualizing module (Tag 0AWE)}).
Hence $f_*\omega_X \to \omega_A$ is surjective whence an isomorphism.
Thus $Rf_*\omega_X = \omega_A$ which by duality implies
$Rf_*\mathcal{O}_X = \mathcal{O}_{\Spec(A)}$.
Whence $H^1(X, \mathcal{O}_X) = 0$ which implies that $A$
defines a rational singularity (see discussion in
Resolution of Surfaces, Section
\href{https://stacks.math.columbia.edu/tag/0AXE}{section bounded (Tag 0AXE)} in particular
Lemmas \href{https://stacks.math.columbia.edu/tag/0B4Q}{lemma regular rational (Tag 0B4Q)} and
\href{https://stacks.math.columbia.edu/tag/0AXF}{lemma exact sequence (Tag 0AXF)}).
If $f$ is minimal, then $E = 0$ because the map
$f^*\omega_A \to \omega_X$ is surjective by
a repeated application of Resolution of Surfaces, Lemma
\href{https://stacks.math.columbia.edu/tag/0B64}{lemma dualizing blow up rational (Tag 0B64)}
and $\omega_A \cong A$ as we've seen above.
\end{proof}


\begin{lemma}
\label{duality-lemma-fundamental-class-lci}
Let $Y$ be a Noetherian scheme. Let $f : X \to Y$ be a
local complete intersection morphism which factors
as an immersion $X \to P$ followed by a proper smooth morphism $P \to Y$.
Let $r$ be the locally constant function on
$X$ such that $\omega_{X/Y} = H^{-r}(f^!\mathcal{O}_Y)$
is the unique nonzero cohomology sheaf of $f^!\mathcal{O}_Y$, see
Lemma \href{https://stacks.math.columbia.edu/tag/0B6V}{lemma lci shriek (Tag 0B6V)}.
Then there is a map
$$
\wedge^r\Omega_{X/Y} \longrightarrow \omega_{X/Y}
$$
which is an isomorphism on the stalk at a point $x$ if and only
if $f$ is smooth at $x$.
\end{lemma}

\begin{proof}
The assumption implies that $X$ is compactifiable over $Y$ hence $f^!$
is defined, see Section \href{https://stacks.math.columbia.edu/tag/0A9Y}{section upper shriek (Tag 0A9Y)}.
Let $j : W \to P$ be an open subscheme such that
$X \to P$ factors through a closed immersion $i : X \to W$.
Moreover, we have $f^! = i^! \circ j^! \circ g^!$ where
$g : P \to Y$ is the given morphism.
We have $g^!\mathcal{O}_Y = \wedge^d\Omega_{P/Y}[d]$ by
Lemma \href{https://stacks.math.columbia.edu/tag/0BRT}{lemma smooth proper (Tag 0BRT)} where $d$ is the locally
constant function giving the relative dimension of $P$ over $Y$.
We have $j^! = j^*$. We have $i^!\mathcal{O}_W = \wedge^c\mathcal{N}[-c]$
where $c$ is the codimension of $X$ in $W$ (a locally constant
function on $X$) and where $\mathcal{N}$ is the normal sheaf of
the Koszul-regular immersion $i$, see Lemma \href{https://stacks.math.columbia.edu/tag/0BR0}{lemma regular immersion (Tag 0BR0)}.
Combining the above we find
$$
f^!\mathcal{O}_Y =
\left(\wedge^c\mathcal{N} \otimes_{\mathcal{O}_X}
\wedge^d\Omega_{P/Y}|_X\right)[d - c]
$$
where we have also used Lemma \href{https://stacks.math.columbia.edu/tag/0B6U}{lemma perfect comparison shriek (Tag 0B6U)}.
Thus $r = d|_X - c$ as locally constant functions on $X$.
The conormal sheaf of $X \to P$ is the module
$\mathcal{I}/\mathcal{I}^2$ where $\mathcal{I} \subset \mathcal{O}_W$
is the ideal sheaf of $i$, see
Morphisms, Section \href{https://stacks.math.columbia.edu/tag/01R1}{section conormal sheaf (Tag 01R1)}.
Consider the canonical exact sequence
$$
\mathcal{I}/\mathcal{I}^2 \to
\Omega_{P/Y}|_X \to \Omega_{X/Y} \to 0
$$
of Morphisms, Lemma \href{https://stacks.math.columbia.edu/tag/01UZ}{lemma differentials relative immersion (Tag 01UZ)}.
We obtain our map by an application of Lemma \href{https://stacks.math.columbia.edu/tag/0E9Y}{lemma determinant (Tag 0E9Y)}.

\medskip\noindent
If $f$ is smooth at $x$, then the map is an isomorphism by an application of
Lemma \href{https://stacks.math.columbia.edu/tag/0E9Y}{lemma determinant (Tag 0E9Y)}
and the fact that $\Omega_{X/Y}$ is locally free at $x$
of rank $r$. Conversely, assume that our map is an isomorphism on stalks
at $x$. Then the lemma shows that $\Omega_{X/Y}$ is free of rank $r$
after replacing $X$ by an open neighbourhood of $x$.
On the other hand, we may also assume that $X = \Spec(A)$ and
$Y = \Spec(R)$ where $A = R[x_1, \ldots, x_n]/(f_1, \ldots, f_m)$
and where $f_1, \ldots, f_m$ is a Koszul regular sequence
(this follows from the definition of local complete intersection morphisms).
Clearly this implies $r = n - m$. We conclude that the rank of the matrix
of partials $\partial f_j/\partial x_i$ in the residue field at $x$ is $m$.
Thus after reordering the variables we may assume
the determinant of $(\partial f_j/\partial x_i)_{1 \leq i, j \leq m}$
is invertible in an open neighbourhood of $x$. It follows
that $R \to A$ is smooth at this point, see for example
Algebra, Example \href{https://stacks.math.columbia.edu/tag/00T8}{example make standard smooth (Tag 00T8)}.
\end{proof}


\begin{lemma}
\label{duality-lemma-fundamental-class-almost-lci}
Let $f : X \to Y$ be a morphism of schemes. Let $r \geq 0$. Assume
\begin{enumerate}
\item $Y$ is Cohen-Macaulay (Properties, Definition
\ref{properties-definition-Cohen-Macaulay}),
\item $f$ factors as $X \to P \to Y$ where the first morphism is
an immersion and the second is smooth and proper,
\item if $x \in X$ and $\dim(\mathcal{O}_{X, x}) \leq 1$,
then $f$ is Koszul at $x$ (More on Morphisms, Definition
\ref{more-morphisms-definition-lci}), and
\item if $\xi$ is a generic point of an irreducible component of $X$, then
we have
$\text{trdeg}_{\kappa(f(\xi))} \kappa(\xi) = r$.
\end{enumerate}
Then with $\omega_{X/Y} = H^{-r}(f^!\mathcal{O}_Y)$ there is a map
$$
\wedge^r\Omega_{X/Y} \longrightarrow \omega_{X/Y}
$$
which is an isomorphism on the locus where $f$ is smooth.
\end{lemma}

\begin{proof}
Let $U \subset X$ be the open subscheme over which $f$ is a
local complete intersection morphism. Since $f$ has relative
dimension $r$ at all generic points by assumption (4) we
see that the locally constant function of
Lemma \ref{duality-lemma-fundamental-class-lci}
is constant with value $r$ and we obtain a map
$$
\wedge^r\Omega_{X/Y}|_U = \wedge^r \Omega_{U/Y}
\longrightarrow
\omega_{U/Y} = \omega_{X/Y}|_U
$$
which is an isomorphism in the smooth points of $f$ (this locus
is contained in $U$ because a smooth morphism is a local complete
intersection morphism). By Lemma \href{https://stacks.math.columbia.edu/tag/0E9V}{lemma shriek over CM (Tag 0E9V)}
and the assumption that $Y$ is Cohen-Macaulay
the module $\omega_{X/Y}$ is $(S_2)$.
Since $U$ contains all the points of codimension $1$ by condition (3)
and using Divisors, Lemma \href{https://stacks.math.columbia.edu/tag/0E9I}{lemma depth 2 hartog (Tag 0E9I)}
we see that $j_*\omega_{U/Y} = \omega_{X/Y}$.
Hence the map over $U$ extends to $X$ and the proof
is complete.
\end{proof}


\begin{lemma}
\label{limits-lemma-glueing-near-closed-point}
Let $S$ be a scheme. Let $s \in S$ be a closed point such that
$U = S \setminus \{s\} \to S$ is quasi-compact. With
$V = \Spec(\mathcal{O}_{S, s}) \setminus \{s\}$ there is
an equivalence of categories
$$
\left\{
\begin{matrix}
X \to S\text{ of finite presentation}
\end{matrix}
\right\}
\longrightarrow
\left\{
\vcenter{
\xymatrix{
X' \ar[d] & Y' \ar[d] \ar[l] \ar[r] & Y \ar[d] \\
U & V \ar[l] \ar[r] & \Spec(\mathcal{O}_{S, s})
}
}
\right\}
$$
where on the right hand side we consider commutative diagrams
whose squares are cartesian and whose vertical arrows are
of finite presentation.
\end{lemma}

\begin{proof}
Let $W \subset S$ be an open neighbourhood of $s$. By
glueing of relative schemes, see
Constructions, Section \href{https://stacks.math.columbia.edu/tag/01LG}{section relative glueing (Tag 01LG)},
the functor
$$
\left\{
\begin{matrix}
X \to S\text{ of finite presentation}
\end{matrix}
\right\}
\longrightarrow
\left\{
\vcenter{
\xymatrix{
X' \ar[d] & Y' \ar[d] \ar[l] \ar[r] & Y \ar[d] \\
U & W \setminus \{s\} \ar[l] \ar[r] & W
}
}
\right\}
$$
is an equivalence of categories. We have
$\mathcal{O}_{S, s} = \colim \mathcal{O}_W(W)$ where
$W$ runs over the affine open neighbourhoods of $s$.
Hence $\Spec(\mathcal{O}_{S, s}) = \lim W$ where $W$
runs over the affine open neighbourhoods of $s$.
Thus the category of schemes of finite presentation
over $\Spec(\mathcal{O}_{S, s})$ is the limit of the
category of schemes of finite presentation over
$W$ where $W$ runs over the affine open neighbourhoods
of $s$, see
Lemma \ref{limits-lemma-descend-finite-presentation}.
For every affine open $s \in W$ we see that $U \cap W$
is quasi-compact as $U \to S$ is quasi-compact.
Hence $V = \lim W \cap U = \lim W \setminus \{s\}$ is a limit of
quasi-compact and quasi-separated schemes (see
Lemma \href{https://stacks.math.columbia.edu/tag/01YX}{lemma directed inverse system has limit (Tag 01YX)}).
Thus also the category of schemes of finite presentation
over $V$ is the limit of the
categories of schemes of finite presentation over
$W \cap U$ where $W$ runs over the affine open neighbourhoods
of $s$. The lemma follows formally from a combination
of these results.
\end{proof}


\begin{lemma}
\label{limits-lemma-descend-finite-presentation}
Let $I$ be a directed set.
Let $(S_i, f_{ii'})$ be an inverse system of schemes over $I$.
Assume
\begin{enumerate}
\item the morphisms $f_{ii'} : S_i \to S_{i'}$ are affine,
\item the schemes $S_i$ are quasi-compact and quasi-separated.
\end{enumerate}
Let $S = \lim_i S_i$. Then we have the following:
\begin{enumerate}
\item For any morphism of finite presentation $X \to S$
there exists an index $i \in I$ and a morphism of finite
presentation $X_i \to S_i$ such that $X \cong X_{i, S}$ as
schemes over $S$.
\item Given an index $i \in I$, schemes
$X_i$, $Y_i$ of finite presentation over $S_i$, and a morphism
$\varphi : X_{i, S} \to Y_{i, S}$ over $S$, there exists an index
$i' \geq i$ and a morphism
$\varphi_{i'} : X_{i, S_{i'}} \to Y_{i, S_{i'}}$
whose base change to $S$ is $\varphi$.
\item Given an index $i \in I$, schemes $X_i$, $Y_i$ of finite presentation
over $S_i$ and a pair of morphisms $\varphi_i, \psi_i : X_i \to Y_i$
whose base changes $\varphi_{i, S} = \psi_{i, S}$ are equal,
there exists an index $i' \geq i$ such that
$\varphi_{i, S_{i'}} = \psi_{i, S_{i'}}$.
\end{enumerate}
In other words, the category of schemes of finite presentation over
$S$ is the colimit over $I$ of the categories of schemes of finite
presentation over $S_i$.
\end{lemma}

\begin{proof}
In case each of the schemes $S_i$ is affine, and we consider
only affine schemes of finite presentation over $S_i$, resp.\ $S$
this lemma is equivalent to
Algebra, Lemma \href{https://stacks.math.columbia.edu/tag/05N9}{lemma colimit category fp algebras (Tag 05N9)}.
We claim that the affine case implies the lemma in general.

\medskip\noindent
Let us prove (3). Suppose given an index $i \in I$, schemes
$X_i$, $Y_i$ of finite presentation over $S_i$ and a pair of morphisms
$\varphi_i, \psi_i : X_i \to Y_i$. Assume that the base changes are
equal: $\varphi_{i, S} = \psi_{i, S}$. We will use the notation
$X_{i'} = X_{i, S_{i'}}$ and $Y_{i'} = Y_{i, S_{i'}}$ for
$i' \geq i$. We also set $X = X_{i, S}$ and $Y = Y_{i, S}$.
Note that according to Lemma \href{https://stacks.math.columbia.edu/tag/01YZ}{lemma scheme over limit (Tag 01YZ)} we have
$X = \lim_{i' \geq i} X_{i'}$ and similarly for $Y$.
Additionally we denote $\varphi_{i'}$ and $\psi_{i'}$
(resp.\ $\varphi$ and $\psi$)
the base change of $\varphi_i$ and $\psi_i$ to $S_{i'}$
(resp.\ $S$). So our assumption means that $\varphi = \psi$.
Since $Y_i$ and $X_i$ are of finite presentation
over $S_i$, and since $S_i$ is quasi-compact and quasi-separated, also
$X_i$ and $Y_i$ are quasi-compact and quasi-separated
(see Morphisms,
Lemma \href{https://stacks.math.columbia.edu/tag/01TY}{lemma finite presentation quasi compact quasi separated (Tag 01TY)}).
Hence we may choose a finite affine open covering
$Y_i = \bigcup V_{j, i}$ such that each $V_{j, i}$ maps into
an affine open of $S_i$. As above, denote $V_{j, i'}$ the inverse
image of $V_{j, i}$ in $Y_{i'}$ and $V_j$ the inverse image in $Y$.
The immersions $V_{j, i'} \to Y_{i'}$ are quasi-compact, and the inverse images
$U_{j, i'} = \varphi_i^{-1}(V_{j, i'})$ and
$U_{j, i'}' = \psi_i^{-1}(V_{j, i'})$
are quasi-compact opens of $X_{i'}$. By assumption the inverse images of
$V_j$ under $\varphi$ and $\psi$ in $X$ are equal.
Hence by Lemma \href{https://stacks.math.columbia.edu/tag/01Z4}{lemma descend opens (Tag 01Z4)}
there exists an index $i' \geq i$ such that
of $U_{j, i'} = U_{j, i'}'$ in $X_{i'}$.
Choose an finite affine open covering
$U_{j, i'} = U_{j, i'}' = \bigcup W_{j, k, i'}$
which induce coverings $U_{j, i''} = U_{j, i''}' = \bigcup W_{j, k, i''}$
for all $i'' \geq i'$.
By the affine case there exists
an index $i''$ such that
$\varphi_{i''}|_{W_{j, k, i''}} = \psi_{i''}|_{W_{j, k, i''}}$
for all $j, k$. Then $i''$ is an index such that
$\varphi_{i''} = \psi_{i''}$ and (3) is proved.

\medskip\noindent
Let us prove (2). Suppose given an index $i \in I$, schemes
$X_i$, $Y_i$ of finite presentation over $S_i$ and a morphism
$\varphi : X_{i, S} \to Y_{i, S}$. We will use the notation
$X_{i'} = X_{i, S_{i'}}$ and $Y_{i'} = Y_{i, S_{i'}}$ for
$i' \geq i$. We also set $X = X_{i, S}$ and $Y = Y_{i, S}$.
Note that according to Lemma \href{https://stacks.math.columbia.edu/tag/01YZ}{lemma scheme over limit (Tag 01YZ)} we have
$X = \lim_{i' \geq i} X_{i'}$ and similarly for $Y$.
Since $Y_i$ and $X_i$ are of finite presentation
over $S_i$, and since $S_i$ is quasi-compact and quasi-separated, also
$X_i$ and $Y_i$ are quasi-compact and quasi-separated
(see Morphisms,
Lemma \href{https://stacks.math.columbia.edu/tag/01TY}{lemma finite presentation quasi compact quasi separated (Tag 01TY)}).
Hence we may choose a finite affine open covering
$Y_i = \bigcup V_{j, i}$ such that each $V_{j, i}$ maps into
an affine open of $S_i$. As above, denote $V_{j, i'}$ the inverse
image of $V_{j, i}$ in $Y_{i'}$ and $V_j$ the inverse image in $Y$.
The immersions $V_j \to Y$ are quasi-compact, and the inverse images
$U_j = \varphi^{-1}(V_j)$ are quasi-compact opens of $X$.
Hence by Lemma \href{https://stacks.math.columbia.edu/tag/01Z4}{lemma descend opens (Tag 01Z4)} there exists an index
$i' \geq i$ and quasi-compact opens $U_{j, i'}$ of $X_{i'}$
whose inverse image in $X$ is $U_j$. Choose an finite affine open covering
$U_{j, i'} = \bigcup W_{j, k, i'}$ which induce affine open
coverings $U_{j, i''} = \bigcup W_{j, k, i''}$
for all $i'' \geq i'$ and an affine open covering
$U_j = \bigcup W_{j, k}$. By the affine case there exists
an index $i''$ and morphisms
$\varphi_{j, k, i''} : W_{j, k, i''} \to V_{j, i''}$
such that
$\varphi|_{W_{j, k}} = \varphi_{j, k, i'', S}$ for all $j, k$.
By part (3) proved above, there is a further index $i''' \geq i''$
such that
$$
\varphi_{j_1, k_1, i'', S_{i'''}}|_{W_{j_1, k_1, i'''} \cap W_{j_2, k_2, i'''}}
=
\varphi_{j_2, k_2, i'', S_{i'''}}|_{W_{j_1, k_1, i'''} \cap W_{j_2, k_2, i'''}}
$$
for all $j_1, j_2, k_1, k_2$. Then $i'''$ is an index such that
there exists a morphism $\varphi_{i'''} : X_{i'''} \to Y_{i'''}$
whose base change to $S$ gives $\varphi$. Hence (2) holds.

\medskip\noindent
Let us prove (1). Suppose given a scheme $X$ of finite presentation
over $S$. Since $X$ is of finite presentation
over $S$, and since $S$ is quasi-compact and quasi-separated, also
$X$ is quasi-compact and quasi-separated
(see Morphisms,
Lemma \href{https://stacks.math.columbia.edu/tag/01TY}{lemma finite presentation quasi compact quasi separated (Tag 01TY)}).
Choose a finite affine open covering $X = \bigcup U_j$
such that each $U_j$ maps into an affine open $V_j \subset S$.
Denote $U_{j_1j_2} = U_{j_1} \cap U_{j_2}$ and
$U_{j_1j_2j_3} = U_{j_1} \cap U_{j_2} \cap U_{j_3}$.
By Lemmas \href{https://stacks.math.columbia.edu/tag/01Z4}{lemma descend opens (Tag 01Z4)} and \href{https://stacks.math.columbia.edu/tag/01Z6}{lemma limit affine (Tag 01Z6)}
we can find an index $i_1$ and affine opens $V_{j, i_1} \subset S_{i_1}$
such that each $V_j$ is the inverse of this in $S$.
Let $V_{j, i}$ be the inverse image of $V_{j, i_1}$ in $S_i$ for
$i \geq i_1$. By the affine case we may find an index $i_2 \geq i_1$ and
affine schemes $U_{j, i_2} \to V_{j, i_2}$ such
that $U_j = S \times_{S_{i_2}} U_{j, i_2}$ is the base change.
Denote $U_{j, i} = S_i \times_{S_{i_2}} U_{j, i_2}$ for $i \geq i_2$.
By Lemma \href{https://stacks.math.columbia.edu/tag/01Z4}{lemma descend opens (Tag 01Z4)} there exists an index
$i_3 \geq i_2$ and open subschemes
$W_{j_1, j_2, i_3} \subset U_{j_1, i_3}$
whose base change to $S$ is equal to $U_{j_1j_2}$.
Denote $W_{j_1, j_2, i} = S_i \times_{S_{i_3}} W_{j_1, j_2, i_3}$
for $i \geq i_3$. By part (2) shown above there exists an index
$i_4 \geq i_3$ and morphisms
$\varphi_{j_1, j_2, i_4} : W_{j_1, j_2, i_4} \to W_{j_2, j_1, i_4}$
whose base change to $S$ gives the identity morphism
$U_{j_1j_2} = U_{j_2j_1}$ for all $j_1, j_2$.
For all $i \geq i_4$ denote
$\varphi_{j_1, j_2, i} = \text{id}_S \times \varphi_{j_1, j_2, i_4}$
the base change. We claim that for some $i_5 \geq i_4$ the system
$((U_{j, i_5})_j, (W_{j_1, j_2, i_5})_{j_1, j_2},
(\varphi_{j_1, j_2, i_5})_{j_1, j_2})$ forms a glueing datum
as in Schemes, Section \href{https://stacks.math.columbia.edu/tag/01JA}{section glueing schemes (Tag 01JA)}.
In order to see this we have to verify that for $i$ large enough
we have
$$
\varphi_{j_1, j_2, i}^{-1}(W_{j_2, j_1, i} \cap W_{j_2, j_3, i})
=
W_{j_1, j_2, i} \cap W_{j_1, j_3, i}
$$
and that for large enough $i$ the cocycle condition holds.
The first condition follows from Lemma \href{https://stacks.math.columbia.edu/tag/01Z4}{lemma descend opens (Tag 01Z4)}
and the fact that $U_{j_2j_1j_3} = U_{j_1j_2j_3}$.
The second from part (3) of the lemma proved above and the fact
that the cocycle condition holds for the maps
$\text{id} : U_{j_1j_2} \to U_{j_2j_1}$.
Ok, so now we can use Schemes, Lemma \href{https://stacks.math.columbia.edu/tag/01JC}{lemma glue schemes (Tag 01JC)}
to glue the system
$((U_{j, i_5})_j, (W_{j_1, j_2, i_5})_{j_1, j_2},
(\varphi_{j_1, j_2, i_5})_{j_1, j_2})$ to get a scheme
$X_{i_5} \to S_{i_5}$. By construction the base change of
$X_{i_5}$ to $S$ is formed by glueing the open affines
$U_j$ along the opens $U_{j_1} \leftarrow U_{j_1j_2} \rightarrow U_{j_2}$.
Hence $S \times_{S_{i_5}} X_{i_5} \cong X$ as desired.
\end{proof}


\begin{lemma}
\label{limits-lemma-descend-closed-immersion-finite-presentation}
Notation and assumptions as in Situation \ref{limits-situation-descent-property}.
If
\begin{enumerate}
\item $f$ is a closed immersion, and
\item $f_0$ is locally of finite type,
\end{enumerate}
then there exists an $i \geq 0$ such that $f_i$ is a closed immersion.
\end{lemma}

\begin{proof}
A closed immersion is affine, see
Morphisms, Lemma \href{https://stacks.math.columbia.edu/tag/01SE}{lemma closed immersion affine (Tag 01SE)}.
Hence by Lemma \href{https://stacks.math.columbia.edu/tag/01ZN}{lemma descend affine finite presentation (Tag 01ZN)} above
after increasing $0$ we may assume that $f_0$ is affine.
By writing $Y_0$ as a finite union of affines we reduce to proving
the result when $X_0$ and $Y_0$ are affine and map
into a common affine $W \subset S_0$. The corresponding algebra
statement is a consequence of
Algebra, Lemma \href{https://stacks.math.columbia.edu/tag/07RH}{lemma colimit surjective (Tag 07RH)}.
\end{proof}


\begin{lemma}[Purity of ramification locus]
\label{lemma-purity-ramification}
\begin{reference}
This result for complex spaces can be found on page 170 of \href{https://stacks.math.columbia.edu/bibliography}{Bibliography: Fischer}.
In general this is \href{https://stacks.math.columbia.edu/bibliography}{Bibliography: Zong (Theorem 2.4)} attributed to Gabber.
\end{reference}
Let $f : X \to Y$ be a morphism of locally Noetherian schemes.
Let $x \in X$ and set $y = f(x)$. Assume
\begin{enumerate}
\item $\mathcal{O}_{X, x}$ is normal of dimension $\geq 1$,
\item $\mathcal{O}_{Y, y}$ is regular,
\item $f$ is locally of finite type, and
\item for specializations $x' \leadsto x$ with
$\dim(\mathcal{O}_{X, x'}) = 1$ our $f$ is \'etale at $x'$.
\end{enumerate}
Then $f$ is \'etale at $x$.
\end{lemma}

\begin{proof}
We will prove the lemma by induction on $d = \dim(\mathcal{O}_{X, x})$.

\medskip\noindent
An uninteresting case is $d = 1$ since in that case the morphism
$f$ is \'etale at $x$ by assumption. Assume $d \geq 2$.

\medskip\noindent
We can base change by $\Spec(\mathcal{O}_{Y, y}) \to Y$
without affecting the conclusion of the lemma, see
Morphisms, Lemma \href{https://stacks.math.columbia.edu/tag/0476}{lemma set points where fibres etale (Tag 0476)}.
Thus we may assume $Y = \Spec(A)$ where $A$ is a regular local
ring and $y$ corresponds to the maximal ideal $\mathfrak m$ of $A$.

\medskip\noindent
Let $x' \leadsto x$ be a specialization with $x' \not = x$.
Then $\mathcal{O}_{X, x'}$ is normal as a localization of
$\mathcal{O}_{X, x}$. If $x'$ is not a generic point of $X$,
then $1 \leq \dim(\mathcal{O}_{X, x'}) < d$ and we conclude that
$f$ is \'etale at $x'$ by induction hypothesis.
Thus we may assume that $f$ is \'etale at all points specializing to
$x$. Since the set of points where $f$ is \'etale is open in $X$
(by definition) we may after replacing $X$ by an open neighbourhood of $x$
assume that $f$ is \'etale away from $\overline{\{x\}}$.
In particular, we see that $f$ is \'etale except at points
lying over the closed point $y \in Y = \Spec(A)$.

\medskip\noindent
Let $X' = X \times_{\Spec(A)} \Spec(A^\wedge)$. Let $x' \in X'$
be the unique point lying over $x$. By the above we see that
$X'$ is \'etale over $\Spec(A^\wedge)$ away from the closed fibre and
hence $X'$ is normal away from the closed fibre. Since $X$ is normal
we conclude that $X'$ is normal by
Resolution of Surfaces, Lemma \href{https://stacks.math.columbia.edu/tag/0BG9}{lemma normalization completion (Tag 0BG9)}.
Then if we can show $X' \to \Spec(A^\wedge)$ is \'etale at $x'$,
then $f$ is \'etale at $x$ (by the aforementioned
Morphisms, Lemma \href{https://stacks.math.columbia.edu/tag/0476}{lemma set points where fibres etale (Tag 0476)}).
Thus we may and do assume $A$ is a regular complete local ring.

\medskip\noindent
The case $d = 2$ now follows from Lemma \ref{lemma-key-purity-ramification}.

\medskip\noindent
Assume $d > 2$. Let $t \in \mathfrak m$, $t \not \in \mathfrak m^2$.
Set $Y_0 = \Spec(A/tA)$ and $X_0 = X \times_Y Y_0$.
Then $X_0 \to Y_0$ is \'etale away from the fibre over the closed point.
Since $d > 2$ we have $\dim(\mathcal{O}_{X_0, x}) = d - 1$ is $\geq 2$.
The normalization $X_0' \to X_0$ is surjective and finite
(as we're working over a complete local ring and such rings are Nagata).
Let $x' \in X_0'$ be a point mapping to $x$. By induction hypothesis the
morphism $X'_0 \to Y$ is \'etale at $x'$. From the inclusions
$\kappa(y) \subset \kappa(x) \subset \kappa(x')$
we conclude that $\kappa(x)$ is finite over $\kappa(y)$.
Hence $x$ is a closed point of the fibre of $X \to Y$
over $y$. But since $x$ is also a generic point of
this fibre, we conclude that $f$ is quasi-finite at $x$
and we reduce to the case of purity of branch locus, see
Lemma \ref{lemma-purity}.
\end{proof}
\end{document}
